\documentclass[11pt,a4paper]{article}

\usepackage[czech]{babel}
\usepackage[utf8]{inputenc}
\usepackage[T1]{fontenc}
\usepackage[a4paper,margin=2.3cm]{geometry}
\usepackage{graphicx}
\usepackage{tikz}
\usetikzlibrary{automata,arrows.meta,positioning,shapes.geometric}
\usepackage{longtable,array,booktabs}
\usepackage{amsmath}
\usepackage{hyperref}
\usepackage{xcolor}
\usepackage{enumitem}
\usepackage{listings}
\usepackage{microtype}

\hypersetup{
  hidelinks,
  pdftitle={Dokumentace projektu IFJ26},
  pdfauthor={Tým IFJ26}
}

\setlength{\parindent}{0pt}
\setlength{\parskip}{0.45em}
\renewcommand{\arraystretch}{1.12}
\lstset{
  basicstyle=\ttfamily\small,
  columns=fullflexible,
  keepspaces=true,
  breaklines=true
}
\newcommand{\placeholder}[1]{\textcolor{red}{\textbf{[DOPLNIT: #1]}}}
\newcommand{\idnt}[1]{\texttt{#1}}


\begin{document}

\begin{titlepage}
\begin{center}
  {\LARGE \textbf{Dokumentace projektu IFJ26}}\\[0.5cm]
  {\large překladač zdrojového jazyka}\\[1.2cm]
\end{center}

\textbf{Varianta:} \placeholder{vv-BVS/TRP-izp}\\[0.25cm]

\begin{table}[h]
\centering
\begin{tabular}{>{\bfseries}p{3.7cm}p{4.2cm}p{2.8cm}p{2.7cm}}
\toprule
Jméno a příjmení & VUT login & Osobní číslo & Podíl bodů \\
\midrule
\textbf{Martin Knotek} & \placeholder{xknotem00} & \placeholder{DOPLNUJTE} & \placeholder{např. 34\,\%} \\
Filip Maxa & \placeholder{VUT login} & \placeholder{DOPLNUJTE} & \placeholder{např. 33\,\%} \\
\placeholder{Jméno třetího člena} & \placeholder{xmiklej00} & \placeholder{DOPLNUJTE} & \placeholder{např. 33\,\%} \\
\bottomrule
\end{tabular}
\end{table}

\textbf{Identifikátory implementovaných rozšíření:}
\placeholder{např. BASE, FUNEXP; uvede se ``-'' při jejich absenci}\\[0.25cm]

\textbf{Odkaz na Git repozitář:}\\
\url{https://github.com/Mett686/IFJ26}\\[0.25cm]

\textbf{Rozdělení bodů:}
\placeholder{doložit součtem 100\,\% a povinně odůvodnit odchylky od rovnoměrného rozdělení}\\
\textbf{Součet bodů:} \placeholder{100\,\%}

\vfill
\begin{center}
  {\small Zdrojový text dokumentace: \texttt{dokumentace.tex}}
\end{center}
\end{titlepage}

\tableofcontents
\newpage

\section{Rozdělení práce v týmu}
\begin{itemize}[leftmargin=*]
  \item \textbf{Martin Knotek -- } \placeholder{syntaktická/sémantická analýza, generování kódu, integrační práce apod.}
  \item \textbf{Filip Maxa -- } \placeholder{lexikální analyzátor, tokenizace, vstupní automat a testování skeneru apod.}
  \item \textbf{\placeholder{Třetí člen} -- } \placeholder{konkrétní moduly a odpovědnost; doplňte před odevzdáním.}
\end{itemize}

\placeholder{Odchylky od rovnoměrného rozdělení bodů zdůvodněte několika konkrétními větami; zaměřte se na rozsah, složitost a důsledky práce, nikoli na spotřebovaný čas.}

\section{Spolupráce, testování a formátování}
\subsection{Spolupráce a code review}
\placeholder{Zde stručně popište větvení, PR, povinná schválení, minimální počet reviewerů a pravidlo pro merge (např. zelené CI a jedno schválení).}

\subsection{Testování}
\placeholder{Uveďte jednotkami vět či odrážek jednotkové testy skeneru, tabulky symbolů a sémantiky, integrační a regresní testy a způsob přípravy negativních vstupů.}

\subsubsection{Použité testovací příkazy}
\begin{lstlisting}
make build
make test
./compiler vstup.ifj26
\end{lstlisting}

\subsection{Styl formátování kódu}
\placeholder{Doplňte odsazení, pojmenování typů a funkcí, délku řádků, použití \texttt{clang-format}/editorconfig či interní konvence.}

\section{Architektura překladače}
\subsection{Přehled komponent}
\begin{figure}[ht]
\centering
\begin{tikzpicture}[
  node distance=7mm and 10mm,
  module/.style={rectangle,rounded corners,draw=black,fill=blue!5,
                 minimum height=11mm,minimum width=31mm,align=center},
  data/.style={rectangle,draw=gray,dashed,fill=yellow!8,
               minimum height=8mm,minimum width=27mm,align=center},
  io/.style={rectangle,rounded corners,draw=black,fill=green!8,
             minimum height=10mm,minimum width=22mm,align=center},
  arrow/.style={-{Latex[length=2mm]},thick}
]
\node[io] (in) {Zdrojový text};
\node[module,below=of in] (scanner) {Lexikální analyzátor\\\texttt{scanner.c/.h}};
\node[data,below=of scanner] (tokens) {Tokiny\\\texttt{token\_t}};
\node[module,below=of tokens] (parser) {Syntaktický analyzátor\\\texttt{parser.c/.h}};
\node[data,below=of parser] (ast) {Abstraktní synt. strom\\\texttt{AST\_Node}};
\node[module,below=of ast] (semantic) {Sémantická analýza\\\texttt{abtDecoding.c/.h}};
\node[data,below=of semantic] (validated) {Validovaný AST + tabulka symbolů};
\node[module,below=of validated] (codegen) {Generátor kódu\\\texttt{codegen.c/.h}};
\node[io,below=of codegen] (out) {IFJcode26};
\draw[arrow] (in) -- (scanner);
\draw[arrow] (scanner) -- node[right]{tokeny} (tokens);
\draw[arrow] (tokens) -- (parser);
\draw[arrow] (parser) -- node[right]{AST} (ast);
\draw[arrow] (ast) -- (semantic);
\draw[arrow] (semantic) -- node[right]{kontrola typů, oborů a zdrojů} (validated);
\draw[arrow] (validated) -- (codegen);
\draw[arrow] (codegen) -- (out);

\node[data,left=of parser] (stack) {Prioritní zásobník\\\texttt{stack.c/.h}};
\draw[arrow] (stack) -- (parser);
\node[data,left=of semantic] (htab) {Rozbalená htab\\\texttt{htab.c/.h}};
\draw[arrow] (htab) -- (semantic);
\end{tikzpicture}
\caption{Blokové schéma toku dat}
\end{figure}

\subsection{Popis toku dat}
\placeholder{Jasně popište vstup a výstup, formáty \texttt{token\_t}, AST, položek tabulky symbolů a cílového IFJcode27. U každé komponenty zdůrazněte odpovědnost a vazby.}

\section{Lexikální analyzátor}
\subsection{Návrh konečného automatu}
Lexikální analyzátor čte vstupní soubor po značkách a vrací jeden \texttt{token\_t}.
Implementace stavového automatu je umístěna ve \texttt{scanner.c}; stavy jsou definovány
ve \texttt{scanner\_private.h}. Indentace je řešena zvláštními stavem a zásobníkem.

\begin{figure}[ht]
\centering
\begin{tikzpicture}[->,>=Stealth,shorten >=1pt,auto,node distance=3.2cm,
                    every state/.style={minimum size=9mm}]
\node[state,initial] (start) {START};
\node[state] (id) [above right=of start] {ID};
\node[state] (int) [right=of start] {CELÉ ČÍSLO};
\node[state] (dbl) [below right=of start] {DES. ČÍSLO};
\node[state] (str) [below=of start] {ŘETĚZEC};
\node[state] (op) [left=of start] {OPERÁTOR};
\node[state] (indent) [above=of start] {ODSAZENÍ};
\node[state,accepting] (token) [below=of str] {PŘEDEJ TOKEN};
\draw (start) edge[bend left=12] node{písmeno/\_} (id);
\draw (start) edge node{číslice} (int);
\draw (start) edge[bend left=12] node{tečka po číslici} (dbl);
\draw (start) edge node{\texttt{\char34}} (str);
\draw (start) edge[bend left=12] node{+-*/<>=!} (op);
\draw (start) edge node{začátek řádku} (indent);
\draw (id) edge[bend left=12] node{klíčové slovo} (token);
\draw (int) edge node{. e E} (dbl);
\draw (str) edge node{uzavírací znak} (token);
\draw (op) edge[bend left=12] node{komplet. lexém} (token);
\draw (indent) edge[bend left=12] node{urovniou/DEDENT} (token);
\draw (int) edge[bend left=12] node{oddělovač} (token);
\draw (dbl) edge[bend left=12] node{oddělovač} (token);
\draw (token) edge[bend right=18] node{další znak} (start);
\end{tikzpicture}
\caption{Zjednodušený přehled stavů lexikálního analyzátoru}
\end{figure}

\placeholder{Nahraďte zjednodušený diagram úplným automatem podle skutečných stavů \texttt{EState}; vyznačte počáteční stav, koncové stavy, přechody a lexikální chyby. Slovně popište odlišnosti od standardního modelu DFA.}

\subsection{Základní lexikální pravidla}
\begin{itemize}
  \item identifikátor: \placeholder{doplňte přesné pravidlo};
  \item celočíselný literál: \placeholder{doplňte přesné pravidlo a rozsah};
  \item desetinný a exponenciální literál: \placeholder{doplňte přesné pravidlo};
  \item řetězec: \placeholder{jednotlivý i vícenásobný řetězec a escape sekvence};
  \item komentáře: \placeholder{formát a pravidla pro blokové/řádkové komentáře};
  \item odchylky v odsazení: \placeholder{kolik mezer znamená úroveň a zpracování EOF po odsazení}.
\end{itemize}

\section{Syntaktický analyzátor}
\subsection{LL gramatika}
\placeholder{Úplnou gramatiku přepište do sloupců níže. Zachovejte terminály, neterminály, epsilon a jednoznačné názvy pravidel. Výchozí neterminál je \texttt{<program>}.}
\[
\begin{array}{rcl}
\langle program\rangle &\to& \texttt{import ifj26.* NEWLINE }\langle definitions\rangle\\
\langle definitions\rangle &\to& \langle definition\rangle\ \langle definitions\rangle\mid\varepsilon\\
\langle \mathrm{TODO}\rangle &\to& \mathrm{TODO}\\
\end{array}
\]
\textit{Poznámka:} Náhradní sada pravidel je založena na souboru
\texttt{SyntaktickyAnalizPravidla.txt}; před odevzdáním ji doplňte podle zdrojové implementace.

\subsection{LL tabulka}
\begin{table}[ht]
\centering
\footnotesize
\setlength{\tabcolsep}{4pt}
\begin{tabular}{l|*{8}{c}}
Neterminál & import & def & @main & ID & if & while & return & \$ \\
\hline
$\langle program\rangle$ & \placeholder{č. pravidla} & & & & & & & \\
$\langle definitions\rangle$ & & \placeholder{} & \placeholder{} & & & & & \\
$\langle statement\rangle$ & & & & \placeholder{} & \placeholder{} & \placeholder{} & \placeholder{} & \\
$\langle \mathrm{TODO}\rangle$ & & & & & & & & \\
\end{tabular}
\caption{LL syntaktická tabulka; doplňte všechna pravidla a neterminály}
\end{table}

\subsection{Precedenční analýza výrazů}
Výrazy jsou analyzované shora dolů v \texttt{parser.c} a pomocí priority zásobníku
vyplněného pomocnými funkcemi ze \texttt{stack.c}. Precedence symbolů a operátorů je
definovaná statickou tabulkou \texttt{table} ve \texttt{parser.h}.

\begin{table}[ht]
\centering
\scriptsize
\setlength{\tabcolsep}{3.2pt}
\begin{tabular}{l|*{16}{c}}
 & + & * & / & = & $>$ & $<$ & $\leq$ & $\geq$ & $\neq$ & - & ( & ) & \$ & ID & \placeholder{U+} & \placeholder{U-} \\
\hline
+ & $>$ & $<$ & $<$ & $>$ & $>$ & $>$ & $>$ & $>$ & $>$ & $>$ & $<$ & $>$ & $>$ & $<$ & $<$ & $<$ \\
* & $>$ & $>$ & $>$ & $>$ & $>$ & $>$ & $>$ & $>$ & $>$ & $>$ & $<$ & $>$ & $>$ & $<$ & $<$ & $<$ \\
\placeholder{doplňte} & & & & & & & & & & & & & & & & \\
\end{tabular}
\caption{Precedenční tabulka implementovaná ve \texttt{parser.h}}
\end{table}
\placeholder{Přidejte úplnou tabulku 16$\times$16 ze zdrojového kódu nebo její čitelnou ekvivalentní formu. Tabulku musí odpovídat stejná priority, asociativita a pravidla pro závorky.}

\section{Sémantická analýza a symbolová tabulka}
\placeholder{Popište kontrolu definic, redefinic, typu proměnných a funkcí, kompatibility typů,
návratových hodnot, oborů platnosti a povinnosti definice @main.}

Implementace využívá rozbalenou hash tabulku \texttt{htab\_t}.
Položka \texttt{htabItem\_t} obsahuje klíč, kontext, mutabilitu, funkční parametry a návratový typ.
\placeholder{Uveďte přesnou významovou a životnostní sémantiku lokálních a globálních položek.}

\section{Generování cílového kódu}
\placeholder{Popište průchod AST, přidělování rámců a proměnných, výrazů, podmínek, cyklů,
volání funkcí a mapování IFJ26 typů na IFJcode26. Uveďte známá omezení implementace.}

\subsection{Mezikód}
\placeholder{Popsat \texttt{TACInstruction}/operandovou reprezentaci, pokud je použita, a jejímu vztahu ke generovanému \texttt{IFJcode26}.}

\section{Stručný popis modulů}
\begin{longtable}{p{3.4cm}p{11.2cm}}
\toprule
Soubory & Úloha a význam \\
\midrule
\texttt{main.c} & \placeholder{Řízení vstupu, alokace symbolové tabulky, volání parseru, sémantické kontroly a generování výstupu.} \\
\texttt{scanner.c/.h}, \texttt{scanner\_private.h} & \placeholder{Stavový automat, lexikální tokeny, vypočet odsazení a řetězcové hodnoty.} \\
\texttt{parser.c/.h} & \placeholder{LL rekurzivní části, precedenční analýza a konstrukce AST.} \\
\texttt{ast.h} & \placeholder{Typy uzlů, operátorů, dat a stromová reprezentace programu.} \\
\texttt{stack.c/.h} & \placeholder{Práce s terminály/neterminály a daty priority zásobníku při analýze výrazů.} \\
\texttt{abtDecoding.c/.h} & \placeholder{Průchod AST a sémantické kontroly, hlášení chyb a typová dedukce.} \\
\texttt{htab.c/.h}, \texttt{htab\_private.h} & \placeholder{Datová struktura tabulky symbolů, hashování, lookup a rozbalování oborů.} \\
\texttt{codegen.c/.h} & \placeholder{Generování instrukcí IFJcode26 z validovaného AST.} \\
\texttt{tac.h} & \placeholder{Definice tříadresných instrukcí/operandů, pokud je tento modul aktivně součástí implementace. Jinak jej explicitně označte jako rozpracovanou či rezervní reprezentaci.} \\
\texttt{SyntaktickyAnalizPravidla.txt} & \placeholder{Ladicí zápis LL pravidel použité jako východisko pro syntaktickou tabulku.} \\
\bottomrule
\end{longtable}

\section{Použití umělé inteligence}
\placeholder{Popsat nástroje a modely AI, účely použití, kontrolní postupy a případný vygenerovaný kód
(např. návrh testů, ladění chyb, úprava dokumentace). Neuvádějte obsahový výstup AI, ale přesný způsob,
jak byl ověřen člověkem a včetně informací vyžadovaných zadáním.}

\section{Známá omezení a testované případy}
\placeholder{Stručně uveďte, co implementace neumí nebo zatím nepodporuje, a co bylo ověřeno.}

\section{Citace převzatých zdrojů}
\begin{thebibliography}{9}
\bibitem{ifj-zadani}
Vysoké učení technické v Brně, \textit{Zadání projektu IFJ26}, akademický rok 2025/2026.
\placeholder{Doplňte přesné vydání a URL použité zadání.}

\bibitem{ifj-teorie}
\placeholder{Úplná bibliografická citace přednáškových či oficiálních materiálů IFJ využitých při návrhu skeneru, parseru a generování kódu.}

\bibitem{htab}
\placeholder{Citace převzaté implementace tabulky symbolů, hashování nebo testů, pokud byla použita.}

\bibitem{ai}
\placeholder{Citace nebo interní URL využitého AI systému; podle požadavků zadání popište původ a účel převzatého návrhu/kódu.}
\end{thebibliography}

\end{document}
